\chapter{积分}\label{ch:b2:integration}

本科一年级分册在闭区间上建立了积分。本章
将其推广到任意区间（广义积分，配以完整的
比较判别法），然后研究\emph{含参
变量的积分} --- 连续性及在积分
号下求导 --- 这是本章唯一
一个在本阶段不加证明地承认的结果。$\Gamma$ 函数充当
贯穿始终的例子，也是通往数学中一半特殊
函数的门径。

\section{任意区间上的积分}

\begin{definition}\label{def:b2:integration:improper}
设 $f$ 为 $\intco{a}{b}$ 上的分段连续函数（$b \in \R$ 或
$+\infty$）。该积分\emph{收敛}\index{广义积分}
当 $\lim_{x \to b^-} \int_a^x f$ 存在时，记
$\int_a^b f$ 为该极限。（类似地，对 $\intoc{a}{b}$ 及
$\intoo{a}{b}$ 上可在某内点处分割 --- 选法无关，
由沙勒关系可知。）称该积分\emph{绝对收敛}，
当 $\int_a^b \abs f$ 收敛时；绝对收敛蕴含
收敛，由柯西准则可得：
\[
\Bigl| \int_x^{y} f \Bigr| \leq \int_x^{y} \abs f
\]
及 $\R$ 的完备性（原函数具有柯西性质）。
详细地说，设 $F(x) = \int_a^x f$ 与 $G(x) = \int_a^x \abs
f$。若 $\int^b\abs f$ 收敛，则 $G$ 在 $b^-$ 处有极限，于是
对任意 $\varepsilon > 0$，存在 $c < b$，使得 $G(y) - G(x)
\leq \varepsilon$ 在 $c \leq x \leq y < b$ 时成立；
该式将此柯西性质传递给 $F$。对任一满足 $x_n
\to b^-$ 的数列，则 $F(x_n)$ 构成实数
的柯西数列，由完备性知其收敛，将两个这样的数列
交错，可知极限与所取数列无关：$F$ 在
$b^-$ 处有极限。
\end{definition}

\begin{theorem}[非负比较判别法]\label{thm:b2:integration:comparison}
设 $f, g \geq 0$ 在 $\intco{a}{b}$ 上分段连续：
\begin{enumerate}
  \item $\int_a^b f$ 收敛，当且仅当原函数 $x \mapsto \int_a^x
        f$ 有界；
  \item $f \leq g$：$\int g$ 收敛蕴含 $\int f$ 收敛；
        发散则反向传递；
  \item $f \sim g$（在 $b$ 处）：两积分敛散性相同；
  \item 参照尺度：在 $+\infty$ 处，$\int^{\infty}
        \frac{\dd t}{t^\alpha}$ 收敛，当且仅当 $\alpha > 1$，且
        $\int^\infty \frac{\dd t}{t(\ln t)^\beta}$ 收敛，当且仅当 $\beta > 1$；
        在有限端点 $b$ 处，$\int^b \frac{\dd t}{(b -
        t)^\alpha}$ 收敛，当且仅当 $\alpha < 1$。
\end{enumerate}
\end{theorem}

\begin{proof}
(1) 原函数 $F(x) = \int_a^x f$ 单调不减（$f \geq
0$）。若其有界，则 $\ell = \sup_{x < b}F$ 有限，且
$F(x) \to \ell$：给定 $\varepsilon > 0$，存在某点，使 $F(x_0) > \ell -
\varepsilon$，由单调性可知 $F(x) \in
\intoc{\ell - \varepsilon}{\ell}$，对 $x_0 \leq x < b$ 成立。若它
无界，则 $F \to +\infty$：发散。

(2) 由 $f \leq g$ 可得 $\int_a^x f \leq \int_a^x g$ 对所有 $x$ 成立；
若 $\int^b g$ 收敛，则右端有界，从而左端亦有界，
由 (1) 即得结论。取逆否即得发散性
反向传递。

(3) $f \sim g$（在 $b$ 处）保证存在 $c < b$，使
\[
\tfrac12\,g(t) \;\leq\; f(t) \;\leq\; 2\,g(t)
\qquad (c \leq t < b) :
\]
由 (2) 双向应用于 $\intco{c}{b}$，两积分
敛散性相同；起始段 $\intcc{a}{c}$ 为
常义积分，不影响结论。

(4) 显式的原函数：对 $\alpha \neq 1$ 与 $\beta \neq 1$，
\[
\int_c^x \frac{\dd t}{t^\alpha}
= \frac{x^{1-\alpha} - c^{1-\alpha}}{1 - \alpha},
\qquad
\int_c^x \frac{\dd t}{t(\ln t)^\beta}
= \frac{(\ln x)^{1-\beta} - (\ln c)^{1-\beta}}{1 - \beta},
\]
在被排除的情形中以对数代替。当 $x \to
+\infty$ 时有界，当且仅当 $\alpha > 1$，对应地为 $\beta > 1$。在
有限端点处，作代换 $u = b - t$，化为
标准尺度 $\int_0 u^{-\alpha}\,\dd u$，当且仅当 $\alpha < 1$ 时有界。
每次均应用 (1)。
\end{proof}

\begin{example}[两个热身例子，完整求解]\label{ex:b2:integration:warmups}
\emph{(a)} $\displaystyle\int_0^1 \ln t\,\dd t$：被积函数
在 $0^+$ 附近无界，但 $\abs{\ln t} = o\bigl(t^{-1/2}\bigr)$
在该处成立（对数增长慢于幂），且 $\int_0 t^{-1/2}$
收敛，即绝对收敛。由分部积分法在
$\intcc{\varepsilon}{1}$ 上计算其值：
\[
\int_\varepsilon^1 \ln t\,\dd t
= \bigl[t\ln t - t\bigr]_\varepsilon^1
= -1 - \varepsilon\ln\varepsilon + \varepsilon
\xrightarrow[\varepsilon\to0^+]{} -1 .
\]
\emph{(b)} $\displaystyle\int_0^\infty \frac{\ln t}{1 +
t^2}\,\dd t$：两端都有瑕点，故在 $1$ 处分割。在 $0$ 附近：
$\abs{\ln t}$ 如 (a) 所述可积；在 $\infty$ 附近：$\frac{\ln
t}{1+t^2} = o(t^{-3/2})$：绝对收敛。
作代换 $t = \frac1u$ 将 $\intoo{0}{1}$ 映为
$\intoo{1}{\infty}$，且
\[
\int_0^1 \frac{\ln t}{1+t^2}\,\dd t
= \int_1^{\infty} \frac{-\ln u}{1 + u^{-2}}\cdot
\frac{\dd u}{u^2}
= -\int_1^\infty \frac{\ln u}{1+u^2}\,\dd u :
\]
两半相互抵消，积分为 $0$。收尾的
体会：$t \mapsto \frac1t$ 下的对称性，抵得上一页
计算 --- 同一技巧已在
\cref{exo:b2:integration:3}中发挥作用。
\end{example}

\begin{example}[同一值，三个积分]\label{ex:b2:integration:onecos}
考察 $I = \displaystyle\int_0^{\infty} \frac{1 - \cos
t}{t^2}\,\dd t$。在 $0$ 处：$1 - \cos t \sim \frac{t^2}2$，故被积函数
可连续延拓，取值 $\frac12$ --- 完全
不存在奇点。在 $\infty$ 处：$0 \leq \frac{1 - \cos t}{t^2}
\leq \frac{2}{t^2}$，故绝对收敛
（见\cref{thm:b2:integration:comparison}）。求值：
在 $\intcc{\varepsilon}{M}$ 上用分部积分，取 $u = 1 - \cos t$，$v' =
t^{-2}$：
\[
\int_\varepsilon^M \frac{1 - \cos t}{t^2}\,\dd t
= \Bigl[-\frac{1 - \cos t}{t}\Bigr]_\varepsilon^M
+ \int_\varepsilon^M \frac{\sin t}{t}\,\dd t .
\]
方括号在两端都为零（$\frac{1 -
\cos\varepsilon}{\varepsilon} \sim \frac\varepsilon2$；
在 $M$ 处分子有界），积分趋于狄利克雷值
$\frac\pi2$（见\cref{exo:b2:integration:10}）：$I = \frac\pi2$。
收尾的体会：利用 $1 - \cos t = 2\sin^2\frac t2$，取 $u =
\frac t2$，
\[
I = \int_0^\infty \frac{2\sin^2 u}{(2u)^2}\;2\,\dd u
= \int_0^\infty \Bigl(\frac{\sin u}{u}\Bigr)^{\!2}\dd u :
\]
三个经典积分 $\int_0^\infty\frac{\sin t}{t}\dd t$，
$\int_0^\infty\bigl(\frac{\sin t}{t}\bigr)^2\dd t$
（见\cref{exo:b2:integration:11}）与 $I$ 的值都是
$\frac\pi2$，经分部积分与换元辗转相传 --- 而
只有第一个是半收敛的：分部积分以
绝对收敛换取了更简单的被积函数。
\end{example}

\begin{example}[半收敛积分]\label{ex:b2:integration:sint}
$\displaystyle\int_1^{\infty} \frac{\sin t}{t}\,\dd t$ 收敛：
用分部积分，
\[
\int_1^x \frac{\sin t}{t}\dd t
= \Bigl[\frac{-\cos t}{t}\Bigr]_1^x - \int_1^x \frac{\cos
t}{t^2}\dd t ,
\]
其中方括号项有极限，最后一个积分
绝对收敛（$\abs{\cos t}/t^2 \leq t^{-2}$）。但原积分并非绝对收敛：
由 $\abs{\sin t} \geq \sin^2 t$，
\[
\int_1^x \frac{\abs{\sin t}}{t}\,\dd t
\;\geq\; \int_1^x \frac{\sin^2t}{t}\,\dd t
= \underbrace{\int_1^x \frac{\dd t}{2t}}_{=\ \frac12\ln x
\ \to\ \infty}
\;-\; \underbrace{\int_1^x \frac{\cos 2t}{2t}\,\dd
t}_{\text{收敛}} ,
\]
其中最后一个积分可由\emph{同样的}
分部积分（方括号中换为 $\sin 2t$）得出：
一个发散项减去一个收敛项，必然发散。故
$\int_1^\infty\frac{\sin t}{t}\dd t$ 收敛，
而不绝对收敛 --- 这是交错级数的积分类比，
其中分部积分扮演的角色
恰如交错判别法。
\end{example}

\section{收敛定理}

\begin{theorem}[控制收敛定理]\label{thm:b2:integration:dominated}
设 $(f_n)$ 为区间 $I$ 上的分段连续函数，
逐点收敛于分段连续函数 $f$，并设存在
一个固定的可积函数 $\varphi \geq 0$（$\int_I \varphi < \infty$），
使得\index{控制收敛定理}
\[
\abs{f_n(t)} \leq \varphi(t) \qquad (t \in I,\ n \in \N).
\]
于是所有 $\int_I f_n$ 与 $\int_I f$ 都绝对收敛，且
\[
\int_I f_n \xrightarrow[n \to \infty]{} \int_I f .
\]
\end{theorem}
\admitted

\begin{remark}
严格的证明属于勒贝格积分理论，将于本科
三年级讲授；但定理本身从此起被频繁使用。
\emph{控制}条件才是关键：仅逐点收敛
不足以保证结论（$f_n = n\,\mathbf{1}_{\intoo{0}{1/n}}$，
滑动凸块：$\int f_n = 1 \not\to 0 = \int f$）。该定理
也对连续参数成立（$f_\lambda$，$\lambda \to
\lambda_0$），由极限的序列刻画即可得到。
\end{remark}

\begin{example}[借助控制收敛求一个高斯极限]\label{ex:b2:integration:wallislimit}
求 $\displaystyle\lim_{n\to\infty} I_n$，其中 $I_n =
\int_0^\infty \Bigl(1 + \frac{t^2}{n}\Bigr)^{\!-n}\dd t$。
逐点地，$(1 + t^2/n)^n \to \eu^{t^2}$（复利极限）
故被积函数趋于 $\eu^{-t^2}$。控制条件：
数列 $n \mapsto (1 + u/n)^n$ 在 $u \geq 0$ 时单调不减
（由算术–几何平均不等式作用于 $n + 1$ 个因子 $1, 1 + \frac un, \dots, 1 +
\frac un$ 得 $(1 + \frac u{n+1})^{n+1} \geq (1 + \frac
un)^n$），因此对 $n \geq 2$：
\[
\Bigl(1 + \frac{t^2}{n}\Bigr)^{\!-n}
\leq \Bigl(1 + \frac{t^2}{2}\Bigr)^{\!-2},
\]
于是得到可积的控制函数（$\sim 4t^{-4}$，于无穷远处）。由
控制收敛定理：
\[
I_n \xrightarrow[n\to\infty]{} \int_0^\infty \eu^{-t^2}\dd t
= \frac{\sqrt\pi}{2}
\]
（高斯积分，见\cref{exo:b2:integration:8}）。收尾的
校验：取换元 $t = \sqrt n\tan\theta$，可算得 $I_n$，
精确地，$I_n = \sqrt n\int_0^{\pi/2}\cos^{2n-2}\theta\,\dd\theta
= \sqrt n\,W_{2n-2}$，而沃利斯渐近式 $W_m \sim
\sqrt{\pi/(2m)}$（见\cref{lem:b2:comparison:wallis}）给出 $\sqrt
n\,W_{2n-2} \to \frac{\sqrt\pi}2$，结果再次相同：本章与上一章的
两大支柱彼此吻合。
\end{example}

\begin{example}[控制收敛，连续参数]
\label{ex:b2:integration:arctanlimit}
求
\[
\lim_{x\to+\infty}\int_0^\infty
\frac{\arctan(xt)}{1+t^2}\,\dd t .
\]
对每个 $t > 0$，
$\arctan(xt) \to \frac\pi2$（$x \to \infty$）；并且
控制：
\[
\Bigl|\frac{\arctan(xt)}{1+t^2}\Bigr|
\leq \frac{\pi/2}{1+t^2},
\qquad\text{可积，与 } x,
\]
对所有 $x$ 成立，且与之无关。由
\cref{thm:b2:integration:dominated}的连续参数形式（序列刻画：
沿每个 $x_n \to \infty$ 检验），
\[
\int_0^\infty\frac{\arctan(xt)}{1+t^2}\,\dd t
\xrightarrow[x\to+\infty]{}
\frac\pi2\int_0^\infty\frac{\dd t}{1+t^2}
= \frac{\pi^2}{4} .
\]
收尾点评：在单独一点 $t = 0$ 处，逐点
极限为 $0$ 而非 $\frac\pi2$，对结果毫无影响 ---
极限函数只通过其积分起作用，这正是该定理
悄然带来的好处之一。
\end{example}

\section{含参变量的积分}

\begin{theorem}[积分号下的连续性]\label{thm:b2:integration:continuity}
设 $f \colon A \times I \to \R$（$A$ 为度量空间，$I$ 为
区间），满足：$t \mapsto f(x, t)$ 对每个
$x$ 分段连续；$x \mapsto f(x, t)$ 对每个 $t$ 连续；并且存在一个
\emph{控制}$\abs{f(x,t)} \leq \varphi(t)$（$\varphi$
在 $I$ 上可积，且与 $x$ 无关）。则
\[
F(x) = \int_I f(x, t)\,\dd t
\]
在 $A$ 上有定义且连续。\index{含参积分}
\end{theorem}

\begin{proof}
有定义性：由控制可得绝对收敛。在
$x_0$ 处的连续性：对任一满足 $x_n \to x_0$ 的数列，函数 $g_n(t) =
f(x_n, t)$ 逐点收敛于 $f(x_0, t)$（由 $x$ 上的连续性）
在固定的控制函数 $\varphi$ 下，由控制收敛定理得
$F(x_n) \to F(x_0)$；由连续性的
序列刻画（见\cref{def:b2:metric:continuity}）得证。
\end{proof}

\begin{theorem}[积分号下的求导]\label{thm:b2:integration:leibnizrule}
设 $f \colon J \times I \to \R$（$J$ 为参数区间），
满足：$t \mapsto f(x,t)$ 在 $I$ 上可积（对每个 $x$）；$x \mapsto
f(x,t)$ 属于 $C^1$ 类（对每个 $t$），偏导数
$\frac{\partial f}{\partial x}$ 关于 $t$ 分段连续，
并且被控制：$\bigl|\frac{\partial f}{\partial x}(x,t)\bigr| \leq
\psi(t)$，其中 $\psi$ 可积。则 $F(x) = \int_I f(x,t)\dd t$
为 $C^1$ 类（在 $J$ 上），且
\[
F'(x) = \int_I \frac{\partial f}{\partial x}(x, t)\,\dd t .
\]
\end{theorem}

\begin{proof}
固定 $x$，取 $h_n \to 0$。差商
\[
\frac{F(x + h_n) - F(x)}{h_n}
= \int_I \frac{f(x + h_n, t) - f(x, t)}{h_n}\,\dd t
\]
的被积函数逐点收敛于 $\frac{\partial f}{\partial
x}(x, t)$，且被 $\psi(t)$ 控制。由中值
不等式（在 $x$ 上应用，$t$ 固定）得：
\[
\Bigl|\frac{f(x + h_n, t) - f(x,t)}{h_n}\Bigr|
\leq \sup_{\xi} \Bigl|\frac{\partial f}{\partial x}(\xi, t)\Bigr|
\leq \psi(t) .
\]
由控制收敛定理，差商的极限为 $\int_I \frac{\partial
f}{\partial x}(x,t)\dd t$，即 $F$ 可导，
导数即如上所述，其连续性由
\cref{thm:b2:integration:continuity}对 $\frac{\partial
f}{\partial x}$ 应用得到。
\end{proof}

\begin{example}[含参积分的公式验证]\label{ex:b2:integration:arctanparam}
设 $F(x) = \displaystyle\int_0^\infty \frac{\dd t}{t^2 + x}$
其中 $x > 0$。对每个 $\intcc{a}{b} \subset \intoo{0}{\infty}$，
被积函数被 $\frac{1}{t^2 + a}$ 控制，该函数可积
且与 $x$ 无关，故 $F$ 连续
（见\cref{thm:b2:integration:continuity}）。此处可用
一个显式值来检验该定理：
\[
F(x) = \Bigl[\frac{1}{\sqrt x}\arctan\frac{t}{\sqrt
x}\Bigr]_0^\infty = \frac{\pi}{2\sqrt x} ,
\]
显然连续。现在对积分号下求导，其中
$x$ 的导数 $-\frac{1}{(t^2+x)^2}$ 在
$\intcc ab$ 上被 $\frac{1}{(t^2+a)^2}$ 控制，且可积，
由\cref{thm:b2:integration:leibnizrule}得
\[
F'(x) = -\int_0^\infty \frac{\dd t}{(t^2 + x)^2}
\qquad\text{而}\qquad
F'(x) = -\frac{\pi}{4}\,x^{-3/2} ,
\]
于是不费力地\emph{算出}一个新积分：
$\int_0^\infty\frac{\dd t}{(t^2+x)^2} = \frac{\pi}{4x^{3/2}}$。
收尾点评：对已知的含参积分求导，
是源源不断得出新公式的“工厂” --- 反复求导可得
$\int_0^\infty\frac{\dd t}{(t^2+1)^n}$，对每个 $n$ 皆成立，
无需任何三角代换。
\end{example}

\begin{method}[反常积分的研究]\label{met:b2:integration:study}
给定 $\int_a^b f$：
\begin{enumerate}
  \item 定位问题：列出 $f$ 无界或区间无穷的端点
        （或内部点），
        并分割区间，
        使每一段恰有一个有问题的端点。
  \item 若 $f$ 在该端附近定号，则
        求一个等价函数，并与
        \cref{thm:b2:integration:comparison}中的参考尺度比较。
  \item 若 $f$ 振荡，先检验 $\abs f$（绝对
        收敛）。若 $\int\abs f$ 发散，则试用分部积分，
        以衰减换取振荡，如\cref{ex:b2:integration:sint}所示；
        下界估计（如
        $\abs{\sin t} \geq \sin^2t$）可识别真正的
        半收敛。
  \item 若要求值而非仅判断敛散，可用分部积分、换元，
        或引入参数（对较简单的积分求导，如
        \cref{ex:b2:integration:arctanparam}和
        \cref{ex:b2:integration:laplace}）。
  \item 对任一计算结果作检验：符号与大致量级须与粗略界相符
        （$\int_0^\infty
        \eu^{-t^2}\dd t \in \intoo{0}{1 + \int_1^\infty
        \eu^{-t}}$，故 $\frac{\sqrt\pi}{2} \approx 0.886$ 是
        合理的）；以及尺度变换下的量纲一致性
        （$t \mapsto \lambda t$ 须使两端
        以相同方式伸缩——这是发现遗漏因子最快的方法）。
\end{enumerate}
\end{method}

\begin{remark}[常见陷阱]
三类反复出现的错误。\emph{(i) 依赖参数的
控制函数：} 控制 $\abs{f(x,t)} \leq \varphi(t)$ 须
对所考虑集合中的 $x$ 一致成立；它通常在
区间 $\intcc ab$ 上成立，但并非整体成立——例如
$\int_0^\infty\eu^{-xt}\dd t$ 不存在对所有 $x > 0$ 均有效的
可积控制函数，但在 $x \geq a > 0$ 上的控制已足以
在整个开半直线上奏效，因为连续性
与导数都是局部性概念。\emph{(ii) 比较有符号的
被积函数：} 比较工具仅适用于\emph{非负}
函数；由 $\abs f \leq g$ 及 $\int g$ 发散，无法
得出任何结论——$\int_1^\infty\frac{\sin t}t\,\dd t$
却收敛，尽管所有与 $\frac1t$ 的比较均告失效。
\emph{(iii) 遗漏一半的问题：}在
$\intoo{0}{\infty}$ 上应始终分别考察两端；
$\int_0^\infty\frac{\dd t}{t}$ 在\emph{两端}均发散，而一个看似收敛的
分解可能悄悄抵消两个无穷大。
稳妥的习惯是逐项核对
\cref{met:b2:integration:study}中的清单。
\end{remark}

\begin{example}[贝特朗边界情形，精确到每一位]\label{ex:b2:integration:bertrandboundary}
$\int^\infty\frac{\dd t}{t(\ln t)^\beta}$ 这一尺度见于
\cref{thm:b2:integration:comparison}中，恰处于幂尺度的边缘，
其边界情形值得完整计算一次。
当 $\beta = 2$ 时：
\[
\int_\eu^{\infty}\frac{\dd t}{t(\ln t)^2}
= \Bigl[-\frac{1}{\ln t}\Bigr]_\eu^{\infty}
= 0 - (-1) = 1 ,
\]
这是一个收敛积分，其值精确而优美；而对
$\beta = 1$，
\[
\int_\eu^{x}\frac{\dd t}{t\ln t}
= \bigl[\ln\ln t\bigr]_\eu^{x} = \ln\ln x
\longrightarrow \infty ,
\]
发散——但发散得如此缓慢，以至于要达到 $\ln\ln x = 10$
需要 $x = \eu^{\eu^{10}} \approx 10^{9566}$。小结：
在“每个幂 $t^{-1-\varepsilon}$
都收敛”与“$t^{-1}$ 发散”之间，存在一个无穷阶梯，
由对数尺度构成，每一级都细化前一级；换元
$u = \ln t$ 把每一级压到前一级上，这正是
贝特朗判别法在高一层上呼应黎曼判别法的
原因。
\end{example}

\begin{remark}[本分册的展望]
本章的工具即将无处不在。控制收敛
定理是下一章各类逼近恒等式的核心驱动力，
无论是滑动核、伯恩斯坦核还是费耶尔核；
积分号下的连续性与求导，
在傅里叶一章中产生了傅里叶系数的微积分，
其中每个 $c_n(f)$ 都是伪装的参数积分。
$\Gamma$ 函数两次登场：一次在
多重积分一章中，那里二重积分终于完整证明了
欧拉的贝塔–伽马公式，另一次在
概率论各章中，那里 $\Gamma$ 型积分用于归一化标准
密度并计算其矩。而半收敛的
$\int\frac{\sin t}{t}$ 在傅里叶一章中以吉布斯
常数的形式再度出现——同一积分
度量了跳跃点处部分和的超调量。
\end{remark}

\begin{definition}[$\Gamma$ 函数]\label{def:b2:integration:gamma}
对 $x > 0$：\index{伽马函数}
\[
\Gamma(x) = \int_0^{\infty} t^{x-1}\,\eu^{-t}\,\dd t ,
\]
该积分在两端收敛（$t^{x-1}$ 在 $0^+$ 处对 $x > 0$ 可积；
在 $\infty$ 处指数衰减）。
\end{definition}

\begin{theorem}\label{thm:b2:integration:gammaprops}
$\Gamma$ 在 $\intoo{0}{+\infty}$ 上连续，且满足
函数方程
\[
\Gamma(x + 1) = x\,\Gamma(x),
\qquad \Gamma(1) = 1,
\qquad\text{从而}\qquad \Gamma(n + 1) = n! ,
\]
且属于 $C^1$ 类（事实上属于 $C^\infty$），并有 $\Gamma'(x) =
\int_0^\infty t^{x-1}\eu^{-t}\ln t\,\dd t$。
\end{theorem}

\begin{proof}
函数方程：在 $\intcc{\varepsilon}{M}$ 上分部积分，
再令两端取极限：$\int t^{x}\eu^{-t} = [-t^x\eu^{-t}] + x\int
t^{x-1}\eu^{-t}$，边界项为零 --- 事实上
$\varepsilon^x\eu^{-\varepsilon} \to 0$，当 $\varepsilon \to
0^+$ 时，因 $x > 0$；又 $M^x\eu^{-M} \to 0$，当 $M \to \infty$
时，因指数函数胜过任何幂次；两个截断积分
都收敛于各自的广义值，
由\cref{def:b2:integration:gamma}中已建立的收敛性即得。$\Gamma(1) = \int
\eu^{-t} = 1$；归纳即得阶乘。

在 $\intcc{a}{b} \subset \intoo{0}{\infty}$ 上证明连续性：取
$t^{x-1}\eu^{-t}$ 的控制函数 $\varphi(t) = (t^{a-1} + t^{b-1})\eu^{-t}$，
它可积，且与 $x \in \intcc{a}{b}$ 无关：
\cref{thm:b2:integration:continuity}在每个这样的区间上适用，
从而在整个半直线上成立。关于可微性，$x$ 的导数
$t^{x-1}\eu^{-t}\ln t$ 在 $\intcc{a}{b}$ 上以 $(t^{a-1}
+ t^{b-1})\eu^{-t}\,\abs{\ln t}$ 为控制，仍可积：
由\cref{thm:b2:integration:leibnizrule}可知；迭代得
到全体各阶导数（每次多出一个 $\ln t$ 因子，无关紧要）。
\end{proof}

\begin{example}[半整数阶乘]\label{ex:b2:integration:halffactorial}
函数方程与 $\Gamma\bigl(\frac12\bigr) =
\sqrt\pi$（只差一个代换，见
\cref{exo:b2:integration:8}：在定义积分中令 $t = u^2$）
可生成全部半整数值：
\[
\Gamma\Bigl(\frac32\Bigr) = \frac12\,\Gamma\Bigl(\frac12\Bigr)
= \frac{\sqrt\pi}{2},
\qquad
\Gamma\Bigl(\frac52\Bigr) = \frac32\cdot\frac{\sqrt\pi}{2}
= \frac{3\sqrt\pi}{4},
\qquad
\Gamma\Bigl(\frac72\Bigr) = \frac{15\sqrt\pi}{8} .
\]
既然 $\Gamma(n+1) = n!$，可以说“$\frac12! =
\frac{\sqrt\pi}{2} \approx 0.886$”：阶乘已被
插值，其曲线降到 $1$ 以下，
位于 $0! = 1$ 与 $1! = 1$ 之间（最小值 $\approx 0.8856$ 在
$x \approx 1.4616$ 处取得，与周末问题
第 I 部分的凸性图像相符）。末了一点：积分
$\int_0^\infty t^{x-1}\eu^{-t}\dd t$ 并不偏爱
整数 --- 阶乘的离散性不过是
计数的偶然，而 $\sqrt\pi$ 正存在于 $1$ 与 $1$ 之间。
\end{example}

\begin{remark}[$\Gamma$ 的后续发展]
本章的周末问题建立了整套欧拉
积分理论，围绕 $\Gamma$：贝塔函数、其分部积分递推式，
作为贝塔值的沃利斯积分，以及高斯
极限公式。多重积分一章以二重积分证明了欧拉的
贝塔–伽马公式对一切自变量成立；概率各章
再次遇到 $\Gamma$，见于最常见密度的归一化，
以及等待时间的矩中。
本科三年级分册在勒贝格积分的基础上重建 $\Gamma$，
证明玻尔–莫勒普唯一性定理，并
将斯特林公式由整数推广到实半直线上，
借助控制收敛定理。
\end{remark}

\begin{example}[通过求导的一个经典计算]\label{ex:b2:integration:laplace}
对 $x \in \R$，记 $F(x) = \int_0^{\infty} \eu^{-t^2}\cos(xt)\,\dd
t$（绝对收敛，由 $\eu^{-t^2}$ 控制）。由
\cref{thm:b2:integration:leibnizrule}（$x$ 的导数
以 $t\,\eu^{-t^2}$ 为控制，可积）：
\[
F'(x) = -\int_0^\infty t\,\eu^{-t^2}\sin(xt)\,\dd t
= \Bigl[\frac{\eu^{-t^2}}{2}\sin(xt)\Bigr]_0^\infty
- \frac x2\int_0^\infty \eu^{-t^2}\cos(xt)\,\dd t
= -\frac x2\,F(x),
\]
（分部积分，取 $u' = t\eu^{-t^2}$）。微分方程 $F' =
-\frac x2 F$ 积分得 $F(x) = F(0)\,\eu^{-x^2/4}$：
高斯型积分自我再现。常数 $F(0) =
\int_0^\infty \eu^{-t^2}\dd t = \frac{\sqrt\pi}{2}$ 的值，
在\cref{exo:b2:integration:8}中算出，并借助二重积分再次
算出，在\cref{ch:b2:multint}中。
\end{example}

\section{练习}

\begin{exercise}[$\star$]\label{exo:b2:integration:1}
判断以下积分的敛散性：
$\displaystyle\int_0^1 \frac{\dd t}{\sqrt{t(1-t)}}$；
$\;\displaystyle\int_1^\infty \frac{\ln t}{t^2}\dd t$；
$\;\displaystyle\int_0^\infty \frac{\dd t}{1 + t^2\sin^2 t}$
\emph{（与 $t =
n\pi$ 附近的发散调和型行为相比较）}。
\end{exercise}

\begin{exercise}[$\star$]\label{exo:b2:integration:2}
计算 $\displaystyle\int_0^\infty t^n \eu^{-\lambda t}\,\dd t$
（$\lambda > 0$）借助 $\Gamma$ 函数，以及
$\displaystyle\int_0^1 (\ln t)^n \dd t$，通过代换 $t =
\eu^{-u}$ 计算。
\end{exercise}

\begin{exercise}[$\star$]\label{exo:b2:integration:3}
证明 $\displaystyle\int_0^{\infty} \frac{\dd t}{(1 +
t^2)(1 + t^x)}$ 对每个 $x \in \R$ 有意义，且
与 $x$ 无关。\emph{（作代换 $t \mapsto \frac1t$，并
将两式取平均。）}它的值是多少？
\end{exercise}

\begin{exercise}[$\star\star$]\label{exo:b2:integration:4}
（有限端点处的贝特朗积分）对哪些 $(\alpha,
\beta)$ 使 $\displaystyle\int_0^{1/2}
\frac{\dd t}{t^\alpha\,\abs{\ln t}^\beta}$ 收敛？
\end{exercise}

\begin{exercise}[$\star\star$]\label{exo:b2:integration:5}
设 $F(x) = \displaystyle\int_0^{\infty} \frac{\eu^{-xt}}{1 +
t^2}\,\dd t$，其中 $x \geq 0$。证明 $F$ 在
$\intco{0}{\infty}$ 上连续，$C^2$ 性质在 $\intoo{0}{\infty}$ 上成立，并在其上满足 $F'' +
F = \frac1x$，且 $F(x) \to 0$（$x \to +\infty$）。
\end{exercise}

\begin{exercise}[$\star\star$]\label{exo:b2:integration:6}
（弗鲁拉尼积分）设 $f$ 在 $\intco{0}{+\infty}$ 上连续，
并有有限极限 $f(\infty)$（在 $+\infty$ 处）。证明：对 $a, b > 0$，
\[
\int_0^{\infty} \frac{f(at) - f(bt)}{t}\,\dd t
= \bigl(f(0) - f(\infty)\bigr)\,\ln\frac ba .
\]
\emph{（在 $\intcc{\varepsilon}{M}$ 上对每一部分作代换，
并将结果重新组合为 $\int_{a\varepsilon}^{b\varepsilon} -
\int_{aM}^{bM}$ 关于 $\frac{f(u)}u\,\dd u$ 的积分，
再利用 $0$ 处的连续性与 $\infty$ 处的极限夹逼。）}计算
$\int_0^\infty \frac{\eu^{-t} - \eu^{-2t}}{t}\dd t$。
\end{exercise}

\begin{exercise}[$\star\star$]\label{exo:b2:integration:7}
证明并计算 $\lim_{n\to\infty} \displaystyle\int_0^n
\Bigl(1 - \frac tn\Bigr)^{\!n} t^{x-1}\,\dd t$，其中 $x > 0$
\emph{（用控制收敛定理，取 $\varphi(t) = \eu^{-t}t^{x-1}$，
利用 $(1 - t/n)^n \leq \eu^{-t}$；极限为 $\Gamma(x)$）}。
\end{exercise}

\begin{exercise}[$\star\star\star$]\label{exo:b2:integration:8}
（用参数技巧计算高斯积分）设 $x \geq 0$，令
\[
G(x) = \Bigl(\int_0^x \eu^{-t^2}\dd t\Bigr)^{\!2},
\qquad
H(x) = \int_0^1 \frac{\eu^{-x^2(1+t^2)}}{1 + t^2}\,\dd t .
\]
证明 $G' + H' = 0$（在积分号下对 $H$ 求导，并
在所得积分中作代换 $u = xt$），由此推出 $G(x) + H(x)
= \frac\pi4$ 对所有 $x$ 成立，并得出
\[
\int_0^{\infty} \eu^{-t^2}\,\dd t = \frac{\sqrt\pi}{2} .
\]
\end{exercise}

\begin{exercise}[$\star\star\star$]\label{exo:b2:integration:9}
证明 $\Gamma$ 为\emph{对数凸}函数，即 $\ln\Gamma$ 在
$\intoo{0}{\infty}$ 上为凸函数。\emph{（对积分形式的柯西--施瓦茨不等式
应用于 $t^{(x+y)/2 - 1}\eu^{-t} = \bigl(t^{x-1}\eu^{-t}\bigr)^{1/2}
\bigl(t^{y-1}\eu^{-t}\bigr)^{1/2}$ 可得
$\Gamma\bigl(\frac{x+y}{2}\bigr)^2 \leq \Gamma(x)\Gamma(y)$；
再结合连续性，并利用\cref{exo:b2:realfun:8}。）}
\end{exercise}

\begin{exercise}[$\star\star\star$]\label{exo:b2:integration:10}
（狄利克雷积分）设 $F(x) = \displaystyle\int_0^{\infty}
\frac{\sin t}{t}\,\eu^{-xt}\,\dd t$，其中 $x > 0$。
\begin{enumerate}
  \item 证明 $F'(x) = -\frac{1}{1 + x^2}$（在
        积分号下求导；计算 $\int_0^\infty \eu^{-xt}\sin t\,\dd
        t$，通过两次分部积分。）
  \item 证明 $F(x) \to 0$（$x \to +\infty$），并由此推出 $F(x) =
        \frac\pi2 - \arctan x$。
  \item 承认 $F$ 在 $0^+$ 处的连续性（阿贝尔型
        定理），求半收敛
        积分的值：
        \[
        \int_0^{\infty} \frac{\sin t}{t}\,\dd t = \frac{\pi}{2}.
        \]
\end{enumerate}
\end{exercise}

\begin{exercise}[$\star\star$]\label{exo:b2:integration:11}
证明 $\displaystyle\int_0^\infty
\Bigl(\frac{\sin t}{t}\Bigr)^{\!2}\dd t$ 收敛，然后用一次
分部积分及\cref{exo:b2:integration:10}计算其值：
\[
\int_0^\infty \Bigl(\frac{\sin t}{t}\Bigr)^{\!2}\dd t
= \frac{\pi}{2} .
\]
（其值与 $\int_0^\infty \frac{\sin t}{t}\dd t$ 相同 --- 但
这一次是绝对收敛。）
\end{exercise}

\begin{exercise}[$\star\star\star$]\label{exo:b2:integration:12}
（高斯积分的尾部）设 $x > 0$，令 $T(x) = \displaystyle
\int_x^\infty \eu^{-t^2}\dd t$。
\begin{enumerate}
  \item 注意到 $\eu^{-t^2} = \frac{1}{-2t}\cdot(-2t\,\eu^{-t^2})$，
        连续两次分部积分，得：
        \[
        T(x) = \eu^{-x^2}\Bigl(\frac{1}{2x} -
        \frac{1}{4x^3}\Bigr) +
        \frac34\int_x^\infty \frac{\eu^{-t^2}}{t^4}\,\dd t .
        \]
  \item 估计余项：$0 \leq \frac34\int_x^\infty
        t^{-4}\eu^{-t^2}\dd t \leq \frac{3}{8x^5}\,\eu^{-x^2}$，
        并由此得如下夹逼：
        \[
        \eu^{-x^2}\Bigl(\frac{1}{2x} - \frac{1}{4x^3}\Bigr)
        \leq T(x) \leq \frac{\eu^{-x^2}}{2x},
        \qquad\text{因此}\qquad
        T(x) \sim \frac{\eu^{-x^2}}{2x} \quad (x \to +\infty).
        \]
  \item 为何由反复分部积分得到的完整交错级数
        对固定的 $x$ 永不收敛？\emph{（比较
        系数 $1\cdot3\cdots(2k-1)$ 的增长与
        幂 $(2x^2)^k$ 的增长。）}
\end{enumerate}
\end{exercise}

\section{问题：欧拉积分 --- 贝塔函数、伽马函数与
高斯极限公式}

\begin{problem}\label{pb:b2:integration:1}
由\cref{def:b2:integration:gamma}给出的 $\Gamma$ 函数，是欧拉积分学的一半；
另一半则是
\emph{贝塔函数}\index{贝塔函数}
\[
B(x, y) = \int_0^1 t^{x-1}(1 - t)^{y-1}\,\dd t .
\]
本题仅用本章的工具 --- 分部积分、换元法、控制收敛 --- 发展 $(\Gamma, B)$ 这一对函数，
并最终得到
\emph{欧拉贝塔–伽马公式}
$B(x,y) = \frac{\Gamma(x)\Gamma(y)}
{\Gamma(x+y)}$（在半整数上成立），以及
\emph{高斯极限公式}\index{高斯极限公式}（针对 $\Gamma$）。
沿途，\cref{lem:b2:comparison:wallis}中的沃利斯积分以贝塔函数值的形式再次出现，
勒让德倍元公式也随之导出。

\medskip
\noindent\textbf{第 I 部分 --- $\Gamma$ 的精细结构。}
\begin{enumerate}
  \item 回想 $\Gamma(x) = \int_0^\infty t^{x-1}\eu^{-t}\dd
        t$ 为何恰在 $x > 0$ 时收敛，并证明
        \[
        \Gamma(x) \sim \frac1x \qquad (x \to 0^+)
        \]
        \emph{（由函数方程及 $\Gamma$ 在
        $1$ 处的连续性）}。
  \item 证明 $\Gamma\bigl(\tfrac12\bigr) = \sqrt\pi$
        \emph{（作代换 $t = u^2$，并用
        \cref{exo:b2:integration:8}）}，并推出
        $\int_\R \eu^{-u^2/2}\dd u = \sqrt{2\pi}$。
  \item 对 $n \in \N$，用归纳法证明：
        \[
        \Gamma\Bigl(n + \frac12\Bigr)
        = \frac{(2n)!}{4^n\,n!}\,\sqrt\pi .
        \]
  \item 说明 $\Gamma''(x) = \int_0^\infty
        t^{x-1}\eu^{-t}(\ln t)^2\dd t > 0$，并由此推出
        $\Gamma$ 严格凸，且在某点
        $x_0 \in \intoo{1}{2}$ 处取唯一极小值，\emph{（由 $\Gamma(1) = \Gamma(2)
        = 1$ 及罗尔定理）}，在 $\intoo{0}{x_0}$ 上递减，且
        在 $\intoo{x_0}{\infty}$ 上递增。
  \item 证明 $\Gamma$ 增长快于任何幂：对每个 $k \in
        \N$，都有 $x^k = o\bigl(\Gamma(x)\bigr)$，当 $x \to +\infty$ 时
        \emph{（用整数夹逼 $x$，并利用 $\Gamma(n+1)
        = n!$ 及第 4 小题的单调性）}。
\end{enumerate}

\medskip
\noindent\textbf{第 II 部分 --- 用分部积分研究贝塔函数。}
\begin{enumerate}[resume]
  \item 证明 $B(x,y)$ 恰在 $x > 0$ 且 $y >
        0$ 时收敛，且 $B(x,y) = B(y,x)$。
  \item 计算 $B(x, 1) = \frac1x$，并
        用分部积分证明，对 $x, y > 0$：
        \[
        B(x, y+1) = \frac{y}{x}\,B(x+1, y) .
        \]
  \item 由分解式 $t^{x-1}(1-t)^{y-1} =
        t^{x}(1-t)^{y-1} + t^{x-1}(1-t)^{y}$，推出 $B(x,y) =
        B(x+1,y) + B(x,y+1)$，并结合第 7 小题，得到
        \emph{递降关系}：
        \[
        B(x, y+1) = \frac{y}{x+y}\,B(x,y),
        \qquad
        B(x+1, y) = \frac{x}{x+y}\,B(x,y) .
        \]
  \item 对整数 $m, n \geq 1$，推出：
        \[
        B(m, n) = \frac{(m-1)!\,(n-1)!}{(m+n-1)!}
        = \frac{1}{(m+n-1)\binom{m+n-2}{m-1}} .
        \]
  \item 证明\emph{一个整数自变量的欧拉公式}：
        对每个 $x > 0$ 及 $n \in \N^*$，
        \[
        B(x, n) = \frac{\Gamma(x)\,\Gamma(n)}{\Gamma(x + n)}
        \]
        \emph{（对 $n$ 作归纳：两边都等于 $\frac1x$（当 $n
        = 1$ 时），且满足相同的递降关系）}。
\end{enumerate}

\medskip
\noindent\textbf{第 III 部分 --- 沃利斯积分作为贝塔函数值。}
\begin{enumerate}[resume]
  \item 作代换 $t = \sin^2\theta$，得到
        三角函数形式
        \[
        B(x, y) = 2\int_0^{\pi/2}
        \sin^{2x-1}\theta\,\cos^{2y-1}\theta\,\dd\theta .
        \]
  \item 推出 $W_n = \frac12\,B\bigl(\frac{n+1}2, \frac12\bigr)$，
        其中沃利斯积分 $W_n = \int_0^{\pi/2}\sin^n
        \theta\,\dd\theta$，并恢复 $W_n =
        \frac{n-1}{n}W_{n-2}$ 这一递推式（见\cref{lem:b2:comparison:wallis}），
        仅由第 8 小题的递降关系即可。
  \item 计算 $B\bigl(\frac12, \frac12\bigr) = 2W_0 = \pi$，并
        将其与 $\Gamma\bigl(\frac12\bigr)^2/\Gamma(1)$ 比较：
        欧拉公式在 $\bigl(\frac12,
        \frac12\bigr)$ 处成立。
  \item 由递推式导出闭形式 $W_{2n} = \frac\pi2\,
        \frac{(2n)!}{4^n(n!)^2}$，并验证
        \[
        B\Bigl(n + \frac12, \frac12\Bigr)
        = \frac{\Gamma\bigl(n + \frac12\bigr)\Gamma\bigl(
        \frac12\bigr)}{\Gamma(n+1)} .
        \]
        借助递降关系用归纳法，可得
        欧拉公式 $B(x,y) =
        \frac{\Gamma(x)\Gamma(y)}{\Gamma(x+y)}$ 在
        $2x$ 与 $2y$ 均为正整数时成立。
  \item 作代换 $u = \frac{t}{1-t}$，得到第三种
        经典形式
        \[
        B(x,y) = \int_0^\infty
        \frac{u^{x-1}}{(1+u)^{x+y}}\,\dd u ,
        \]
        并直接验证 $x = y = \frac12$ 的情形（取 $u = v^2$
        可化为 $\int_0^\infty\frac{2\,\dd v}{1+v^2}$）。
\end{enumerate}

\medskip
\noindent\textbf{第 IV 部分 --- 高斯极限公式。}
\begin{enumerate}[resume]
  \item 设 $x > 0$ 且 $n \in \N^*$，
        用 $n$ 次连续分部积分证明：
        \[
        \int_0^n \Bigl(1 - \frac tn\Bigr)^{\!n} t^{x-1}\,\dd t
        = \frac{n!\;n^x}{x(x+1)\cdots(x+n)} .
        \]
  \item 由 \cref{exo:b2:integration:7}（控制收敛定理）
        得到\emph{高斯极限公式}：
        \[
        \Gamma(x) = \lim_{n\to\infty}
        \frac{n!\;n^x}{x(x+1)\cdots(x+n)}
        \qquad (x > 0).
        \]
  \item 取对数，证明：当 $x > 0$ 时，
        \[
        \ln\Gamma(x) = -\ln x - \gamma x +
        \sum_{k=1}^{\infty}\Bigl(\frac xk -
        \ln\Bigl(1 + \frac xk\Bigr)\Bigr),
        \]
        其中 $\gamma$ 为欧拉常数
        （见\cref{ex:b2:comparison:harmonic}）；
        证明该级数收敛\emph{（一般项
        为 $\sim \frac{x^2}{2k^2}$）}。
  \item 在 $x = \frac12$ 处应用高斯公式，并利用中心二项式系数的
        渐近式 $\binom{2n}{n} \sim
        \frac{4^n}{\sqrt{\pi n}}$
        （见\cref{ex:b2:comparison:centralbinomial}）以重新算出
        $\Gamma\bigl(\frac12\bigr) = \sqrt\pi$：斯特林
        常数与高斯积分是同一个数的
        两种面貌。
  \item 验证高斯公式能重新推出函数
        方程：由下列精确恒等式
        \[
        \frac{n!\,n^{x+1}}{(x+1)\cdots(x+n+1)}
        = \frac{n!\,n^{x}}{x(x+1)\cdots(x+n)}\cdot
        \frac{n\,x}{x+n+1},
        \]
        由此再次得到 $\Gamma(x+1) = x\,\Gamma(x)$。（高斯
        公式已完全确定 $\Gamma$；本科三年级分册
        证明了更精细的玻尔–莫勒鲁普定理：
        函数方程与对数凸性已足以
        唯一确定 $\Gamma$。）
\end{enumerate}

\medskip
\noindent\textbf{第 V 部分 --- 红利。}
\begin{enumerate}[resume]
  \item 设 $a > 0$，证明 $\int_0^\infty \eu^{-t^a}\dd t =
        \Gamma\bigl(1 + \frac1a\bigr)$，并求当
        $a \to +\infty$ 时的极限，用控制收敛定理\emph{（逐点极限
        为 $\mathbf 1_{t < 1}$；取 $1$ 作为 $\intoc{0}{1}$ 上的控制函数，
        其余部分取 $\eu^{-t^2}$ 作为控制函数（$a \geq
        2$ 时）}。并用
        $\Gamma$ 的连续性检验结果。
  \item 设 $n \geq 1$，证明
        \[
        \int_0^1 \frac{\dd t}{\sqrt{1 - t^n}}
        = \frac1n\,B\Bigl(\frac1n, \frac12\Bigr),
        \]
        并求得 $2$（$n = 1$ 时）和 $\frac\pi2$（$n
        = 2$ 时）的值。（当 $n = 4$ 时，这是\emph{双纽线
        常数}，它没有初等封闭形式；其
        来历属于椭圆积分理论。）
  \item （矩）设 $x > 0$ 且 $k \in \N$，证明
        \[
        \frac{1}{\Gamma(x)}\int_0^\infty
        t^{k}\,t^{x-1}\eu^{-t}\,\dd t
        = \frac{\Gamma(x+k)}{\Gamma(x)}
        = x(x+1)\cdots(x+k-1),
        \]
        称为上升阶乘；验证 $x = 1$ 时得 $k!$。（在
        概率论各章中，这是标准等待时间
        密度的 $k$ 阶矩。）
  \item 证明贝塔恒等式，对\emph{所有} $x > 0$ 成立：
        \[
        B(x, x) = 2^{1-2x}\,B\Bigl(x, \frac12\Bigr)
        \]
        \emph{（作代换 $t = \frac{1+s}2$，利用
        关于 $s$ 的对称性，再令 $s = \sqrt v$）}。由此推得，当
        $2x \in \N^*$ 时，\emph{勒让德倍元公式}
        \[
        \Gamma(x)\,\Gamma\Bigl(x + \frac12\Bigr)
        = 2^{1-2x}\,\sqrt\pi\;\Gamma(2x),
        \]
        并由第 3 小题直接验证 $x = n$ 的情形。（对
        一般的 $x$ 则由同一恒等式推出，前提是欧拉公式
        对所有自变量均已成立 --- 其证明见
        多重积分一章中的二重积分
        证明。）
  \item 综合。每项用一句话回答：(i) 分部积分在何处承担了
        第二部分的全部内容；(ii) 控制收敛定理在何处进入
        第四、第五部分；(iii) 哪些渐近性输入
        是从比较一章引入的；(iv) 目前已证明的
        是欧拉公式的哪一部分，
        即 $B(x,y) = \Gamma(x)\Gamma(y)/\Gamma(x+y)$，以及
        还有哪些留待二重积分来解决。
\end{enumerate}
\end{problem}
